% Figura 8.2: Proyectos multifase con resolución progresiva del riesgo (Caso farmacéutico / apertura piloto UGR)
\begin{tikzpicture}[
    font=\small\sffamily,
    fase/.style={
        rectangle, rounded corners=4pt, draw=#1, fill=#1!8,
        line width=1.1pt, minimum width=3.5cm, minimum height=1.35cm,
        align=center, font=\footnotesize\sffamily, text=dfinslate
    },
    hoja/.style={
        rectangle, rounded corners=3pt, draw=#1, fill=#1!10,
        line width=0.9pt, minimum width=4.0cm, minimum height=1.05cm,
        align=center, font=\scriptsize\sffamily, text=dfinslate
    },
    lbl/.style={font=\scriptsize\sffamily, fill=white, inner sep=1.5pt, text=#1}
]

% Nodo t=0: Fase I Piloto / I+D
\node[fase=dfincopper] (t0) at (0, 0) {%
    \textbf{\color{dfincopper}Fase I: Piloto / I+D ($t=0$)}\\[2pt]
    Desembolso inicial $-I_0$\\
    Tasa riesgo alto ($k_1$)
};

% Bifurcación en t=1 (resolución incertidumbre técnica)
\node[fase=dfinemerald] (exito) at (5.4, 1.75) {%
    \textbf{\color{dfinemerald}Éxito técnico ($t=1$)}\\[2pt]
    Probabilidad $p$\\
    Invierte $-I_1$ (expansión)
};

\node[hoja=dfinburg] (fracaso) at (5.4, -1.75) {%
    \textbf{\color{dfinburg}Fracaso / Abandono ($t=1$)}\\[2pt]
    Probabilidad $1-p$\\
    Flujo posterior $= 0$ (no invierte $I_1$)
};

% Fase II Explotación estable t=2..T
\node[hoja=dfinnavy] (explot) at (10.8, 1.75) {%
    \textbf{\color{dfinnavy}Fase II: Explotación ($t=2\dots T$)}\\[2pt]
    Flujos operativos $FCL_t$\\
    Descuento de régimen normal ($k_2 < k_1$)
};

% Conexiones con probabilidades y explicación de valoración inversa
\draw[->, >=Latex, line width=1.2pt, color=dfinemerald] (t0.east) to[out=22, in=180]
    node[lbl=dfinemerald, above=3pt] {$p$ (éxito)} (exito.west);

\draw[->, >=Latex, line width=1.2pt, color=dfinburg] (t0.east) to[out=-22, in=180]
    node[lbl=dfinburg, below=3pt] {$1-p$ (fracaso)} (fracaso.west);

\draw[->, >=Latex, line width=1.2pt, color=dfinnavy] (exito.east) -- (explot.west)
    node[midway, lbl=dfinnavy, above=2pt] {Despliegue};

% Fórmula de inducción hacia atrás debajo
\node[draw=dfinslate!35, fill=dfinnavy!5, rounded corners=3pt, inner sep=4pt, font=\scriptsize\sffamily, text=dfinnavy]
    at (5.4, -3.1) {$\longleftarrow$ Valoración por inducción hacia atrás: $\text{VAN}_0 = -I_0 + \dfrac{p \cdot \max\bigl(0,\ -I_1 + \text{VA}_1(\text{Fase II a } k_2)\bigr)}{1 + k_1}$};

\end{tikzpicture}
